\chapter{Hyperlinks and OSC-8}
\label{ch:hyperlinks}

\begin{aside}
Click a URL in your terminal and a browser opens. The feature is
ten years old and still feels magical. The mechanism is one
escape sequence.
\end{aside}

\section{The OSC-8 escape}

A hyperlink wraps a stretch of cells with two escapes:

\begin{lstlisting}
ESC ] 8 ; ; https://example.com BEL
clickable text
ESC ] 8 ; ; BEL
\end{lstlisting}

The terminal renders the inner text normally (no underline, no
colour change) but stores the URL as metadata on those cells.
Click (or Ctrl-click on some terminals) opens the URL.

\section{Per-cell links}

In \maya{}'s packed cell, a link slot points into a side table
of URLs. Many cells can share one entry (a 50-character link
takes 50 cells, one URL slot). The slot is 14 bits, so up to
16,383 distinct URLs per frame, which is generous.

The diff treats the link slot as part of the cell: a cell whose
URL changed but whose text didn't still emits.

\section{Element wrapper}

\begin{lstlisting}
auto link = hyperlink("https://example.com",
                     text("docs"));
\end{lstlisting}

The wrapper applies the URL to every cell its child paints. It
does not add visual styling --- you can underline if you want
the user to see it:

\begin{lstlisting}
auto link = hyperlink("https://example.com",
                     text("docs") | underline)
          | fg<theme::accent>;
\end{lstlisting}

\section{Detection at runtime}

OSC-8 is silently ignored by terminals that don't support it,
so emitting is always safe. The text shows; clicking does
nothing. \maya{} does not bother with capability detection.

\section{Multi-line links}

A link can span multiple lines. The terminal joins consecutive
cells with the same URL into one clickable region. Wrap a
multi-line block in one \code{hyperlink} call and it works.

\section{IDs for grouping}

OSC-8 supports an ID parameter: cells with the same ID share a
hover/click target. This is useful when the same logical link
appears in multiple places (e.g., a file path repeated in a
search-results view):

\begin{lstlisting}
ESC ] 8 ; id=42 ; file:///etc/passwd BEL
\end{lstlisting}

Hovering one instance highlights all instances with the same
ID. \maya{} accepts an optional \code{id} argument to
\code{hyperlink}.

\section{File and special schemes}

\code{file://}, \code{mailto:}, \code{ssh:}, \code{tel:},
\code{magnet:} --- terminals delegate to the OS to handle the
scheme. \maya{} does no special-casing; you emit the URL and
let the click handler do its job.

\section{Security}

A URL the user sees is not always the URL the link points to:

\begin{lstlisting}
hyperlink("https://malicious.example",
         text("https://google.com"));
\end{lstlisting}

The cell shows ``https://google.com'' but clicks open the
malicious site. Some terminals (iTerm2) show a tooltip with
the real URL on hover; many don't.

\maya{} does not validate or guard against this. The
application must decide whether to trust its sources (your own
formatted strings are fine; user-supplied URLs need care).

\begin{warning}
Echoing untrusted URLs is a phishing vector. If your app
prints URLs from external data (a chat client, an RSS reader),
either strip the OSC-8 metadata or display the visible text
identical to the URL.
\end{warning}

\section{Performance}

OSC-8 is on the wire as plain ASCII. Each link adds 30--80
bytes of escape per emission. Frequent re-emissions (animated
panels with hovers) get expensive.

\maya{}'s style cache (see Chapter~\ref{ch:diffloop}) tracks the
last-emitted link and skips redundant resets. Sequential cells
with the same URL emit one open, one close, regardless of
length.

\section{Terminal support}

\begin{itemize}[leftmargin=*]
  \item \textbf{Supported:} iTerm2, Wezterm, Kitty, foot,
    Ghostty, recent xterm (with flag), modern gnome-terminal.
  \item \textbf{Not supported:} Apple Terminal, old xterm,
    tmux without passthrough.
\end{itemize}

\section{tmux passthrough}

tmux strips OSC-8 by default. Enable passthrough with:

\begin{lstlisting}
set -g allow-passthrough on
\end{lstlisting}

Then links emitted inside a pane reach the outer terminal.
Without this, links are silently dropped at the tmux boundary.

\section{Beyond URLs: app actions}

Some terminals offer the convention of \code{file:///path:line}
to jump to a file at a specific line. iTerm2 and Wezterm
support this. \maya{} treats it as an opaque URL; the
terminal does the parsing.

You can use this to wire ``open file at line'' actions
into your app's output, useful for compiler-error displays.

\section{Recap}

\begin{recap}
\begin{itemize}[leftmargin=*]
  \item OSC-8 wraps cells with metadata; clicking opens.
  \item \maya{} stores the link in the packed cell.
  \item No visual styling unless you add it.
  \item IDs group multi-instance links.
  \item Be careful echoing untrusted URLs.
  \item tmux needs passthrough configured.
\end{itemize}
\end{recap}

\section*{Workshop}

\begin{enumerate}[leftmargin=*]
  \item Print a list of URLs in your app, each as a clickable
    link. Confirm clicking opens your browser.
  \item Build a search-result list where every line is a
    clickable \code{file:///path:line} link.
  \item Stress test: 1000 unique links on screen. Measure
    redraw size.
\end{enumerate}
